%In the past three years,
High capacity supervised learning algorithms, such as deep
convolutional neural networks \cite{lecun1998gradient,AlexNet},
%have enabled significant progress 
have led to a discontinuity in visual recognition over the past four
years and continue to push the state-of-the-art performance (e.g. \cite{MSRA,wu2015deep,GoogleNorm}). These models
usually have millions of parameters, resulting in two consequences.
On the one hand, the many degrees of freedom of the models allow for extremely impressive description power; they can learn to represent complex functions and transformations automatically from the
data. On the other hand, to search for the optimal settings for a
large number of parameters, these data-hungry algorithms require
a massive amount of training data with humainitiallced labels
\cite{deng2009imagenet,fei2007learning,xiao2010sun,zhou2014learning}.

Although there has been remarkable progress in improving deep
learning algorithms (e.g. \cite{MSRA,GoogLeNet}) and developing high
performance training systems (e.g. \cite{wu2015deep}), advancements
are lacking in dataset construction.  The ImageNet dataset
\cite{deng2009imagenet}, used by most of these algorithms, is
7 years old and has been heavily over-fitted \cite{BaiduScandal}.  The 
recently released Places \cite{zhou2014learning} dataset is not much larger.
The number of parameters in
many deep models now exceeds the number of images in these datasets.
%stretching the performance under the risk of the lack of bounded
%constraints \cite{Intriguing,Fooled,AdversarialExamples}.  % Clarify...
While the models are getting deeper
(e.g. \cite{VGGnet,GoogLeNet}) and the accessible computation power is
increasing, the size of the datasets for training and evaluation is
not increasing by much, but rather lagging behind and hindering further
progress in large-scale visual recognition.

Moreover, the density of examples in current datasets
is quite low.  Although they have several million images,
they are spread across many categories, and so there
are not very many images in each category (e.g., 1000 in ImageNet). As a result,
deep networks trained on them often learn features that are noisy
and/or unstable~\cite{Intriguing,Fooled}. To address this problem,
researchers have used techniques that augment the training sets with
perturbations of the original images~\cite{wu2015deep}.  Those methods
add to the stability and generalizability to trained models, but without
increasing the actual density of novel examples.  We propose an
alternative: an image dataset with high density. %but few categories. - why would we advertise few categories?
We aim for a dataset with $\sim10^6$ images per category, which
is around 10 times denser than
PLACES~\cite{zhou2014learning} and 100 times denser than
ImageNet~\cite{deng2009imagenet}.

Although the amount of available image data on the Internet is
increasing constantly, it is nontrivial to build a supervised dataset that
dense because of the high costs of manual labeling.  Clever methods
have been proposed to improve the efficiency of human-in-the-loop
annotation (e.g., \cite{branson2010visual,deng2014scalable,russakovsky2015best,vijayanarasimhan2008multi,wah2011multclass}).
However, manual effort is still the bottleneck -- the
construction of ImageNet and Places both required more than a year of
human effort via Amazon Mechanical Turk (AMT), the largest
crowd-sourcing platform available on the Internet. If we desire a
$N$-times bigger/denser dataset, it will require more than $N$ years of
human annotation. This will not scale quickly enough to support
advancements in deep visual learning. Clearly, we must
introduce a certain amount of automation into this process in order to
let data annotation maintain pace with the growth of deep models.

%However,
%we need to be extremely cautious in ensuring the accuracy of labels
%obtained (semi-)automatically. In addition, humans make mistakes
%quite often, especially when performing repetitive and tedious jobs
%such as data labeling.
% figure to point out errors in places annotation.
% Thus, a systematic way to check human labeling accuracy is a
% necessity. %hopefully even to surpass human accuracy by combining
%automatic algorithms and human annotation. - Not sure %about this

In this paper, we introduce an integrated framework using deep
learning with humans in the loop
%combining deep learning and AMT crowd sourcing 
to annotate a large-scale image dataset. The key new idea is a labeling propagation
system that automatically amplifies manual human effort.  We study strategies for
selecting images for people to label, interfaces for acquiring labels quickly,
procedures for verifying the labels acquired, and methods for amplifying the labels by
propagating them to other images in the dataset.  We have used the system to
construct a large scale image database, ``LSUN'', with 10 million labeled images
in 10 scene categories and 59 million labeled images in 20 object categories.\footnote{We are continuing to collect and
annotate data through our platform at a rate of 5 million images and 2 categories
per week}.

One of the main challenges for a semi-automatic approach
is achieving high precision in the final labeled dataset.  
Our procedure uses statistical tests to ensure labeling quality,
providing more than 90\% precision on average according to
verification tests.  This is slightly lower than manual annotation,
but still good enough to train
high performance classification models.  During experiments with
popular deep learning models, we observe
substantial classifier performance gains when using our larger training
set.  Although the LSUN training labels may contain some noise, it
seems that training on a larger, noisy dataset 
produces better models than training on smaller, noise-free datasets.

% Training vs. test set precision?

%objects too
%    LSUN for scene classification with more than 10 millions images.


